\documentclass[11pt,a4paper]{article}
\usepackage[margin=1in]{geometry}
\usepackage[T1]{fontenc}
\usepackage[utf8]{inputenc}
\usepackage{lmodern}
\usepackage{array}
\usepackage{longtable}
\usepackage{hyperref}
\hypersetup{
  colorlinks=true,
  linkcolor=black,
  urlcolor=blue
}

\title{ME403 Student Simulation Pack Guide}
\author{ME403 Introduction to Robotics}
\date{}

\begin{document}
\maketitle

\section{Purpose}
This document explains how to use the released \texttt{ME403\_LabFiles} package
as a student. It clarifies what is included, what is intentionally excluded, and
how simulation assets are prepared at runtime.

\section{What Is Included}
\begin{itemize}
  \item Lab starter folders under \texttt{labs/}
  \item Vendored SDK and robot description sources under \texttt{vendor/}
  \item Full simulation runtime under \texttt{simulation/runtime/}
  \item Optional simulation demos under \texttt{simulation/student/}
  \item Optional robot setup scripts under \texttt{robots/}
  \item Bootstrap and asset helpers under \texttt{scripts/}
  \item Setup and troubleshooting docs under \texttt{docs/}
\end{itemize}

\section{What Is Not Included by Default}
\begin{itemize}
  \item Instructor/private solution material
  \item Internal course publishing helpers
  \item Build artifacts and temporary caches
\end{itemize}

\section{Quickstart}
Run from package root:
\begin{verbatim}
cd ME403_LabFiles
python3 -m pip install -r requirements/base.txt
python3 scripts/bootstrap.py --simulation
cd labs/lab01_fk
python3 interface.py
\end{verbatim}

\section{Simulation Architecture (Student View)}
\subsection*{Upstream Assets vs Runtime Assets}
The \texttt{vendor/} folder stores upstream robot descriptions and SDK sources.
When URDF backends run, \texttt{simulation/runtime/urdf\_loader.py} prepares
runtime copies for simulator compatibility.

\subsection*{Preparation and Cache Pipeline}
The runtime preparation step can include:
\begin{itemize}
  \item Mesh format conversion for backend compatibility
  \item URI/path rewriting so files resolve correctly
  \item Simulator-specific adjustments and hidden parity chain injection
  \item Cached output placement under \texttt{\$DOBOT\_SIM\_CACHE}
\end{itemize}

Default cache path is:
\begin{verbatim}
~/.cache/dobot_sim
\end{verbatim}

\section{Environment Variables}
\begin{longtable}{|p{0.34\linewidth}|p{0.58\linewidth}|}
\hline
\textbf{Variable} & \textbf{Meaning} \\
\hline
\texttt{DOBOT\_SIMULATION} & Set to \texttt{1} to force simulation mode (default for labs). \\
\hline
\texttt{DOBOT\_SIM\_BACKEND} & Choose \texttt{mujoco} (default) or \texttt{pybullet}. \\
\hline
\texttt{DOBOT\_VIZ} & Set to \texttt{0} for headless mode (no GUI windows). \\
\hline
\texttt{DOBOT\_ROBOT\_TYPE} & Choose \texttt{magician} or \texttt{mg400}. \\
\hline
\texttt{DOBOT\_EE} & Magician TCP mode: \texttt{none} (default, bare flange),
  \texttt{motor}, or \texttt{suction}. Controls which TCP offset the sim uses.
  Must match the real robot's \texttt{set\_tool\_offset()} setting for FK parity. \\
\hline
\texttt{DOBOT\_SIM\_CACHE} & Override cache directory for prepared assets. \\
\hline
\texttt{DOBOT\_MAGICIAN\_URDF\_PATH} & Override Magician URDF source file/folder path. \\
\hline
\texttt{DOBOT\_MG400\_URDF\_PATH} & Override MG400 URDF source file/folder path. \\
\hline
\end{longtable}

\section{Backend Choice}
\begin{itemize}
  \item Use \texttt{mujoco} for the default course simulation workflow.
  \item Use \texttt{pybullet} when MuJoCo graphics/drivers are unavailable.
  \item Use headless mode (\texttt{DOBOT\_VIZ=0}) for remote sessions.
\end{itemize}

\section{Troubleshooting Quick Chart}
\begin{longtable}{|p{0.3\linewidth}|p{0.62\linewidth}|}
\hline
\textbf{Symptom} & \textbf{Action} \\
\hline
Missing URDF/meshes &
Run \texttt{python3 scripts/fetch\_assets.py --all} from package root. \\
\hline
First run is very slow &
Expected behavior: initial fetch + mesh prep + cache generation. \\
\hline
MG400 URDF generation error &
Install \texttt{xacro}, then rerun MG400 asset fetch. \\
\hline
GUI or display issues &
Try \texttt{DOBOT\_VIZ=0}; on Linux test \texttt{MUJOCO\_GL=egl}. \\
\hline
Stale cache behavior &
Delete \texttt{\$DOBOT\_SIM\_CACHE} (or \texttt{\~/.cache/dobot\_sim}) and rerun. \\
\hline
\end{longtable}

\section{Matching Sim to Real-Robot EE Mode (Magician)}

The simulation TCP mode (\texttt{DOBOT\_EE}) and the real robot's stored
offset must agree for FK/IK errors to reflect kinematics rather than an
offset mismatch.

\begin{center}
\begin{tabular}{|l|l|l|}
\hline
\textbf{Sim \texttt{DOBOT\_EE}} & \textbf{TCP offset (mm)} & \textbf{Real robot Python command} \\
\hline
\texttt{none} (default) & (0, 0, 0) & \texttt{set\_tool\_offset(bot, 0, 0, 0)} \\
\hline
\texttt{motor} & (60, 0, 0) & \texttt{set\_tool\_offset(bot, 60, 0, 0)} \\
\hline
\texttt{suction} & (60, 0, $-$70) & \texttt{set\_tool\_offset(bot, 60, 0, -70)} --- physical cup tip \\
\hline
\end{tabular}
\end{center}

\texttt{MAGICIAN\_TOOL\_OFFSETS} dict in \texttt{utils.py} provides shortcuts:
\begin{verbatim}
set_tool_offset(bot, *MAGICIAN_TOOL_OFFSETS["none"])    # FK labs default
\end{verbatim}

For MG400 labs, no TCP offset switching is needed --- the MG400 always uses
its bare wrist reference point in ME403 labs.

\section{Recommended Student Workflow}
\begin{enumerate}
  \item Start in simulation and validate your lab code.
  \item Switch robot type/backend only when needed.
  \item Move to hardware only after simulation results are stable.
  \item For Magician: home once in DobotStudio (v1.9.4), then use
        \texttt{set\_tool\_offset()} from Python to set the TCP mode for each session.
  \item Verify \texttt{DOBOT\_EE} (sim) matches \texttt{get\_tool\_offset(bot)} (real)
        before collecting FK/IK data.
\end{enumerate}

\section{Further Reading}
\begin{itemize}
  \item \texttt{README.md} (root quickstart)
  \item \texttt{docs/simulation.md}
  \item \texttt{docs/troubleshooting.md}
\end{itemize}

\end{document}
